% GENERATED FILE. Edit canonical lessons, then run npm run generate:books.
% canonical-source-hash: fnv1a64:a31fcdb80e4d53da
% canonical-lessons: PA-C208-go-return, PA-C208-pia, PA-C208-pia-review, PA-C208-had-to-go, PA-C208-payoff

\chapter{I Had to Go}
\label{ch:pa-a2-pia}
\hlchaptermodality{208}

\begin{chapteropening}
\textbf{By the end of this chapter:} \emph{I can recognise and use one bounded masculine singular past-necessity form.}

Recognise and produce one bounded I-had-to-go frame through listening, speaking, reading, and writing.
\end{chapteropening}

\section[\texorpdfstring{\pa{ਮੈਨੂੰ} · \pa{ਜਾਣਾ} (mainu; jaana)}{ਮੈਨੂੰ · ਜਾਣਾ (mainu; jaana)}]{Return to to me and to go}

\label{lesson:PA-C208-go-return}



\begin{quote}
\emph{Hear:} \emph{mainu; jaana}; YOU SAY: \textbf{\pa{ਮੈਨੂੰ} · \pa{ਜਾਣਾ}}
\end{quote}

\begin{grammarlens}[title={Grammar Lens}]
Keep the familiar “to me” form and the familiar infinitive “to go” ready. The next lesson will own one new auxiliary base and one exact past form; this opener does not reveal either one.
\end{grammarlens}

\subsection*{Your turn}
\begin{itemize}
\raggedright
  \item \emph{Hear:} \emph{to me}; YOU SAY: \textbf{\pa{ਮੈਨੂੰ}}
  \item \emph{Hear:} \emph{to go}; YOU SAY: \textbf{\pa{ਜਾਣਾ}}
  \item Say the two familiar items once, then stop.
\end{itemize}

\begin{tcolorbox}[breakable,colback=teal!4,colframe=teal!35!black,title={Before you move on}]
Produce “to me” and “to go” from meaning. This opener owns no new form.
\end{tcolorbox}
\section[\texorpdfstring{\pa{ਪੈਣਾ} · \pa{ਪਿਆ} (paina · piaa)}{ਪੈਣਾ · ਪਿਆ (paina · piaa)}]{One bounded past-necessity form}

\label{lesson:PA-C208-pia}



\begin{quote}
Retrieve familiar \textbf{\pa{ਮੈਨੂੰ}}, “to me,” and \textbf{\pa{ਜਾਣਾ}}, “to go.” Then listen for one new auxiliary base and its one past form.
\end{quote}

\begin{grammarlens}[title={Grammar Lens: direct owner}]
\begin{quote}
\textbf{\pa{ਪੈਣਾ}} — \emph{paina} — to have to
\end{quote}

\begin{quote}
\textbf{\pa{ਪਿਆ}} — \emph{piaa} — had to; masculine singular perfective form
\end{quote}

Hear and say the two whole forms before reading them. Keep \textbf{\pa{ਪਿਆ}} whole. This chapter uses it only after one familiar infinitive in one past-necessity frame. It does not open the wider agreement system.
\end{grammarlens}

\begin{cousinweb}
Punjabi University Patiala's reference grammar describes the \textbf{-\pa{ਣਾ} \pa{ਪੈ}-} construction as “have to,” says that the person uses \textbf{\pa{ਨੂੰ}}, and gives both present or future and past examples. AQA's official Panjabi grammar lists the infinitive, past and perfective forms, and auxiliary \textbf{\pa{ਪੈ}}. Sources: \href{https://pt.learnpunjabi.org/assets/A\%20Reference\%20Grammar\_Final.pdf}{Punjabi University reference grammar} and \href{https://www.aqa.org.uk/subjects/panjabi/a-level/panjabi-7682/specification/subject-content/grammar}{AQA Panjabi grammar}.

Those sources support this modern construction. Keep the exact form whole; do not invent a separate etymology for it. This chapter teaches only the masculine singular perfective in one “had to” frame, not a present or future obligation, a feminine or plural form, a full agreement paradigm, or any other sense of the auxiliary.
\end{cousinweb}

\subsection*{Your turn}
\begin{itemize}
\raggedright
  \item \emph{Hear:} \emph{paina}; YOU SAY: \textbf{\pa{ਪੈਣਾ}}
  \item \emph{Hear:} \emph{piaa, had to}; YOU SAY: \textbf{\pa{ਪਿਆ}}
  \item \emph{Read it:} \textbf{\pa{ਪੈਣਾ} · \pa{ਪਿਆ}}
  \item \emph{Write it:} copy \textbf{\pa{ਪਿਆ}} once, checking the whole form
\end{itemize}

\begin{tcolorbox}[breakable,colback=teal!4,colframe=teal!35!black,title={Before you move on}]
Produce the auxiliary base, then its one masculine singular perfective “had to” form.
\end{tcolorbox}
\section[\texorpdfstring{\pa{ਪਿਆ} (piaa)}{ਪਿਆ (piaa)}]{Retrieve the new past-necessity form on its own}

\label{lesson:PA-C208-pia-review}



\begin{quote}
\emph{Hear:} \emph{piaa}; YOU SAY: \textbf{\pa{ਪਿਆ}}
\end{quote}

\subsection*{Your turn}
\begin{itemize}
\raggedright
  \item \emph{Hear:} \emph{had to, masculine singular perfective}; YOU SAY: \textbf{\pa{ਪਿਆ}}
  \item \emph{Read it:} \textbf{\pa{ਪਿਆ}}
  \item \emph{Write it:} cover and write \textbf{\pa{ਪਿਆ}}
\end{itemize}

\begin{grammarlens}[title={Grammar Lens}]
Keep one whole form and one bounded past-necessity meaning here. The familiar words that complete its frame return only after this separate retrieval.
\end{grammarlens}

\begin{tcolorbox}[breakable,colback=teal!4,colframe=teal!35!black,title={Before you move on}]
Say and then write the whole form once from memory.
\end{tcolorbox}
\section[\texorpdfstring{\pa{ਮੈਨੂੰ} \pa{ਜਾਣਾ} \pa{ਪਿਆ} (mainu jaana piaa)}{ਮੈਨੂੰ ਜਾਣਾ ਪਿਆ (mainu jaana piaa)}]{Put the new form in one past-necessity frame}

\label{lesson:PA-C208-had-to-go}



\begin{quote}
Retrieve familiar “to me” and “to go,” then retrieve the new “had to” form.
\end{quote}

\begin{grammarlens}[title={Grammar Lens}]
\begin{quote}
\textbf{\pa{ਮੈਨੂੰ} \pa{ਜਾਣਾ} \pa{ਪਿਆ}} — \emph{mainu jaana piaa} — I had to go
\end{quote}

For this one frame, begin with familiar \textbf{\pa{ਮੈਨੂੰ}}, keep familiar infinitive \textbf{\pa{ਜਾਣਾ}} in the middle, and place new whole form \textbf{\pa{ਪਿਆ}} last. Nothing else in the necessity or agreement system is added.
\end{grammarlens}

\subsection*{Your turn}
\begin{itemize}
\raggedright
  \item \emph{Hear:} \emph{mainu jaana piaa}; YOU SAY: \textbf{\pa{ਮੈਨੂੰ} \pa{ਜਾਣਾ} \pa{ਪਿਆ}}
  \item \emph{Read it:} \textbf{\pa{ਮੈਨੂੰ} \pa{ਜਾਣਾ} \pa{ਪਿਆ}}
  \item \emph{Write it:} cover and write \textbf{\pa{ਮੈਨੂੰ} \pa{ਜਾਣਾ} \pa{ਪਿਆ}}
\end{itemize}

\begin{tcolorbox}[breakable,colback=teal!4,colframe=teal!35!black,title={Before you move on}]
Produce the complete “I had to go” frame.
\end{tcolorbox}
\section[\texorpdfstring{\pa{ਮੈਨੂੰ} \pa{ਜਾਣਾ} \pa{ਪਿਆ} (mainu jaana piaa)}{ਮੈਨੂੰ ਜਾਣਾ ਪਿਆ (mainu jaana piaa)}]{A bounded I-had-to-go payoff}

\label{lesson:PA-C208-payoff}



\begin{quote}
Retrieve the new past-necessity form alone before rebuilding its familiar frame.
\end{quote}

\subsection*{Your turn: four skills}
\begin{itemize}
\raggedright
  \item \emph{Hear:} \emph{mainu jaana piaa}; YOU SAY: \textbf{\pa{ਮੈਨੂੰ} \pa{ਜਾਣਾ} \pa{ਪਿਆ}}
  \item \emph{Read it:} \textbf{\pa{ਮੈਨੂੰ} \pa{ਜਾਣਾ} \pa{ਪਿਆ}}
  \item \emph{Write it:} cover the line and write \textbf{\pa{ਮੈਨੂੰ} \pa{ਜਾਣਾ} \pa{ਪਿਆ}}
\end{itemize}

\begin{tcolorbox}[breakable,colback=teal!4,colframe=teal!35!black,title={Before you move on}]
Say the new form, then produce the complete “I had to go” frame.

This is the thirty-first prerequisite tranche exposed by Punjabi A2 mock drafting. It does not unblock Punjabi A2 mock 1 by itself. It is not an A2 mock, not A2 readiness or attainment, and not external or human validation.
\end{tcolorbox}
